% Part: computability
% Chapter: computability-theory
% Section: fixed-point-thm

\documentclass[../../../include/open-logic-section]{subfiles}

\begin{document}

\olfileid{cmp}{thy}{fix}
\olsection{مقرر نقطے دا قضیہ}

آؤ رُکّن دے مسئلے
نوں فیر ویکھئیے۔
عارضی علامت وجوں
$\gn{\cfind{x}(y)}$
نوں $\tuple{x, y}$
لئی لکھو؛ ایہنوں
$\cfind{x}(y)$ دی
قدر دے اک «ناں»
وجوں سمجھو۔ ایس
علامت نال اسی اپنے اوہناں ثبوتاں وچوں اک
ہور لفظاں وچ بیان کر سکدے آں جیہڑے
وکھاؤندے نیں کہ رُکّن دا مسئلہ ناقابلِ فیصلہ اے۔

سوال: کی اک قابلِ حساب
فنکشن $h$ موجود اے
جیہدی ہِٹھلی خاصیت
ہووے؟ ہر $x$ تے
$y$ لئی:
\[
h(\gn{\cfind{x}(y)}) =
\begin{cases}
1 & \text{جے $\cfind{x}(y) \fdefined$} \\
0 & \text{ہور صورت وچ۔}
\end{cases}
\]

جواب: نہیں؛ ورنہ
ہِٹھاں دتا جزوی فنکشن
\[
g(x) \simeq
\begin{cases}
0 & \text{جے $h(\gn{\cfind{x}(x)}) = 0$} \\
\text{غیر تعریف شدہ} & \text{ہور صورت وچ}
\end{cases}
\]
قابلِ حساب ہووے گا
تے اوس دا کوئی
اشاریہ~$e$ ہووے گا۔
پر فیر:
\[
\cfind{e}(e) \simeq
\begin{cases}
0 & \text{جے $h(\gn{\cfind{e}(e)}) = 0$} \\
\text{غیر تعریف شدہ} & \text{ہور صورت وچ،}
\end{cases}
\]
ایس صورت وچ
$\cfind{e}(e)$
اُدوں تے صرف
اُدوں تعریف شدہ
اے جدوں ایہ
تعریف شدہ نہیں؛
ایہ تضاد اے۔

ہن $\cfind{e}$ والی
مساوات ویکھو۔
اک لحاظ نال ایتھے
اپنے آپ دا حوالہ
اے: اسی $\cfind{e}(e)$
دی قدر نوں اک خاص
طریقے نال
$\gn{\cfind{e}(e)}$
اُتے منحصر بنایا
اے۔ مقرر نقطے دا
قضیہ آکھدا اے کہ
عام طور اُتے اسی
ایہ {\em کر سکدے}
آں؛ صرف تضاد ثابت
کرن لئی نہیں۔

\olref{lem:fixed-equiv}
مقرر نقطے دا قضیہ
بیان کرن دے دو
ہم ارز طریقے دیندا
اے۔ منطقی لحاظ
نال ایہ بیان دونوں
سچ ہون کارن
ہم ارز ہوندے نیں؛
پر ساڈی مراد ایہ
اے کہ ہر اک دوجے
توں سدھا نکلدا اے،
ایس لئی دونوں اکّو قضیے نوں
بیان کرن دے وکھو وکھرے طریقے نیں۔

\begin{lem}
\ollabel{lem:fixed-equiv}
ہِٹھلے بیان ہم ارز نیں:
\begin{enumerate}
\item ہر جزوی قابلِ
  حساب فنکشن $g(x,y)$
  لئی اک اشاریہ~$e$
  اے ایداں کہ ہر~$y$
  لئی:
\[
\cfind{e}(y) \simeq g(e,y).
\]
\item ہر قابلِ حساب
  فنکشن~$f(x)$ لئی
  اک اشاریہ~$e$ اے
  ایداں کہ ہر~$y$
  لئی:
\[
\cfind{e}(y) \simeq \cfind{f(e)}(y).
\]
\end{enumerate}
\end{lem}

\begin{proof}
$(1) \Rightarrow (2)$:
$f$ دتا ہووے تاں
$g$ نوں $g(x,y) \simeq
\fn{Un}(f(x),y)$
دی مدد نال تعریف کرو۔
(1) نال ایداں دا
اشاریہ~$e$ لو کہ
ہر~$y$ لئی:
\begin{align*}
\cfind{e}(y) & = \fn{Un}(f(e),y) \\
& = \cfind{f(e)}(y).
\end{align*}

$(2) \Rightarrow (1)$:
$g$ دتا ہووے تاں
$s$-$m$-$n$ قضیہ
ورت کے $f$ لو
ایداں کہ ہر $x$
تے~$y$ لئی
$\cfind{f(x)}(y) \simeq g(x,y)$۔
(2) نال ایداں دا
اشاریہ~$e$ لو کہ:
\begin{align*}
\cfind{e}(y) & = \cfind{f(e)}(y) \\
& = g(e,y).
\end{align*}
ایس نال ثبوت
مکدا اے۔
\end{proof}

\begin{explain}
بیان (1) ثابت کرن توں پہلاں، جیہدے توں
(2) وی نکلدا اے، سوچو ایہ کِنّا حیران کن اے۔
$e$ نوں کمپیوٹر
پروگرام سمجھو۔
بیان (1) آکھدا
اے کہ کوئی وی
جزوی قابلِ حساب
$g(x,y)$ دتا
ہووے تاں تسی
اک ایداں دا
کمپیوٹر پروگرام
$e$ لبھ سکدے
او جیڑا
$g_e(y) \simeq g(e,y)$
حساب کرے۔
یعنی تسی ایداں
دا پروگرام لبھ
سکدے او جیڑا اوس فنکشن دی قدر
کڈھدا اے جیہڑا ایس پروگرام دا ای حوالہ دیندا اے۔
\end{explain}

\begin{thm}
\olref{lem:fixed-equiv}
دے دونوں بیان
سچ نیں۔ خاص طور
اُتے، ہر جزوی
قابلِ حساب فنکشن
$g(x,y)$ لئی اک
اشاریہ~$e$ اے
ایداں کہ ہر~$y$
لئی:
\[
\cfind{e}(y) \simeq g(e,y).
\]
\end{thm}

\begin{proof}
درکار اجزا اُتّے
رُکّن دے مسئلے
دی بحث وچ پہلے
ای موجود نیں۔
$\fn{diag}(x)$
اک قابلِ حساب
فنکشن ہووے جیڑا
ہر $x$ لئی
$f_x(y) \simeq \cfind{x}(x,y)$
دا اشاریہ دیندا
اے، یعنی:
\[
\cfind{\fn{diag}(x)}(y) \simeq \cfind{x}(x,y).
\]
$\fn{diag}$ نوں
ایداں دا فنکشن
سمجھو جیڑا دو
آرگومنٹاں والے
فنکشن دے پروگرام
نوں اک آرگومنٹ
والے فنکشن دے
پروگرام وچ بدلدا
اے۔ ایس لئی اصل فنکشن دے پہلے آرگومنٹ
وجوں اوہدا اپنا پروگرام، یعنی اوہدا اشاریہ،
مقرر کردے آں۔ $\fn{diag}$
دی رسمی تعریف
ایداں بن سکدی
اے: پہلے $s$
نوں تعریف کرو:
\[
s(x,y) \simeq \fn{Un}^2(x,x,y),
\]
جتھے $\fn{Un}^2$
تِن آرگومنٹاں
والا فنکشن اے
جو دو آرگومنٹاں
والے جزوی قابلِ
حساب فنکشناں
لئی ہمہ گیر اے۔
فیر $s$-$m$-$n$
قضیے نال بنیادی
بازگشتی فنکشن
$\fn{diag}$ لبھدا
اے جیہڑا ایہ
شرط پوری کردا
اے:
\[
\cfind{\fn{diag}(x)}(y) \simeq s(x,y).
\]

ہن فنکشن $l$
نوں ایداں
تعریف کرو:
\[
l(x,y) \simeq g(\fn{diag}(x),y).
\]
تے $\gn{l}$،
$l$ دا اک
اشاریہ ہووے۔
آخر $e = \fn{diag}(\gn{l})$
من لو۔ فیر
ہر $y$ لئی:
\begin{align*}
\cfind{e}(y) & \simeq \cfind{\fn{diag}(\gn{l})}(y) \\
& \simeq \cfind{\gn{l}}(\gn{l}, y) \\
& \simeq l(\gn{l}, y) \\
& \simeq g(\fn{diag}(\gn{l}),y) \\
& \simeq g(e, y),
\end{align*}
جیویں چاہیدا سی۔
\end{proof}

\begin{explain}
ایتھے ہو کی رہیا
اے؟ من لو تہاڈا
کم اک ایداں دا
کمپیوٹر پروگرام
لکھنا اے جیڑا
اپنا متن آپ چھاپے۔
ہور من لو تسی
ایہو جیہی پروگرامنگ
زبان ورتدے او جیہدی
لڑیاں والے فنکشناں
دی لائبریری بھرپور
تے انوکھی اے۔
خاص کر، فرض کرو
ایس زبان وچ
اک فنکشن $\fn{diag}$
اے جیڑا ایداں کم
کردا اے: انپٹ
لڑی~$s$ دتی ہووے
تاں $\fn{diag}$ اوس~$s$ وچ «x» نشان
دیاں ساریاں تھاواں لبھدا اے تے ہر وار
اوس دی تھاں اصل لڑی دا اقتباسی
نشانیاں وچ بند روپ پا دیندا اے۔ مثلاً،
ہِٹھلی لڑی انپٹ
ہووے:
\begin{quote}
\begin{verbatim}
hello x world
\end{verbatim}
\end{quote}
تاں فنکشن ہِٹھلا
نتیجہ دیندا اے:
\begin{quote}
\begin{verbatim}
hello 'hello x world' world
\end{verbatim}
\end{quote}
ایس صورت وچ
مطلوبہ پروگرام
لکھنا سوکھا اے؛
تسی جانچ سکدے
او کہ:
\begin{quote}
\begin{verbatim}
print(diag('print(diag(x))'))
\end{verbatim}
\end{quote}
ایہ کم کردا اے۔
C++ تے Java
ورگیاں عام
پروگرامنگ زباناں
وچ وی اوہی خیال
کم کردا اے،
بھاویں اوہدی
بناوٹ ہور اوکھی
ہووے۔

مقرر نقطے دے
قضیے دے ثبوت
توں اسی ہن
صرف دو کدم
دور آں۔
من لو پرنٹ
فنکشن دی اک
صورت $\fn{print}(x,y)$
لڑی $x$ تے
اک ہور عددی
آرگومنٹ $y$
لیندی اے، تے
لڑی $x$ نوں
$y$ واری
چھاپدی اے۔
تاں ہِٹھلا
«پروگرام»:
\begin{quote}
\begin{verbatim}
getinput(y);
print(diag('getinput(y); print(diag(x), y)'), y)
\end{verbatim}
\end{quote}
انپٹ $y$ اُتے
اپنا متن $y$
واری چھاپدا
اے۔ $\fn{getinput}$---
$\fn{print}$---
$\fn{diag}$ دے
ڈھانچے دی تھاں
کوئی وی فنکشن
$g(x,y)$ رکھو
تاں ملدا اے:
\begin{quote}
\begin{verbatim}
g(diag('g(diag(x), y)'), y)
\end{verbatim}
\end{quote}
ایہ پروگرام
انپٹ $y$
اُتے $g$
نوں ایس پروگرام دے اپنے متن
تے $y$ اُتے
چلاوندا اے۔
«اقتباسی نشانیاں
وچ بند کرن»
نوں «ایس دا
اشاریہ ورتن»
وجوں سمجھو
تاں اُتّے
والا ثبوت
ملدا اے۔

فی الحال جے
تسی ایس ثبوت
نوں رسمی چال
یا کالا جادو
سمجھنا چاہو
تاں کوئی گل
نہیں۔ پر تہانوں اُتّے دتی دلیل دیاں
تفصیلاں آپ دوبارہ بنا سکنا چاہیدا اے۔
جدوں اسی
نامکملیت دے
قضیے، تے
اوہناں نال
متعلق «مقرر
نقطے دا قضیہ»،
ثابت کراں گے
تاں اسی ایس گل نوں سمجھن دے ہور
طریقیاں اُتے گل کراں گے کہ ایہ کم کیوں کردا اے۔
\end{explain}

\begin{tagblock}{lambda}
\begin{digress}
ایہو خیال
«مقرر نقطہ»
کمبینیٹر وی
دیندا اے۔
من لو تہاڈے
کول لیمبڈا اصطلاح
$g$ اے، تے
تسی ہور اصطلاح
$k$ چاہندے
او جیہدی
خاصیت ایہ
ہووے کہ $k$، $\beta$ دے لحاظ نال
$gk$ دے ہم ارز ہووے۔
ساڈے علامتی
قاعدیاں مطابق
ایہ اصطلاحاں
تعریف کرو:
\[
\fn{diag}(x) = xx
\]
تے
\[
l(x) = g(\fn{diag}(x))
\]
ہور لفظاں وچ،
$l$ اصطلاح
$\lambd[x][g(xx)]$
اے۔ $k$ نوں
$ll$ اصطلاح
من لو۔ فیر:
\begin{align*}
k & = (\lambd[x][g(xx)])(\lambd[x][g(xx)]) \\
& \red  g((\lambd[x][g(xx)])(\lambd[x][g(xx)])) \\
& = gk.
\end{align*}
جے ایہ لو:
\[
Y = \lambd[g][((\lambd[x][g(xx)])(\lambd[x][g(xx)]))]
\]
تاں $Yg$ تے
$g(Yg)$ اک
سانجھی اصطلاح
تک گھٹدے نیں؛
ایس لئی $Yg \equiv_\beta
g(Yg)$۔ ایہنوں
«کری دا کمبینیٹر»
آکھدے نیں۔
جے ایس دی
تھاں ایہ لو:
\[
Y = (\lambd[xg][g(xxg)])(\lambd[xg][g(xxg)])
\]
تاں واقعی $Yg$،
$g(Yg)$ ول
گھٹدا اے، جیڑی
ہور مضبوط گل
اے۔ $Y$ دی
ایس دوجی
صورت نوں
«ٹیورنگ دا
کمبینیٹر»
آکھدے نیں۔
\end{digress}
\end{tagblock}

\end{document}
